\section{Interpreting a published benchmark comparison}
\label{sec:comparison-theory}

\subsection{What same-dataset means}
A model comparison on a shared set of FASTA records answers a narrower
question than a paired model comparison on identical individual test
predictions. The Feb2020 site reports a 10-fold CV mean, standard deviation,
and model recipe, but not the exact assignment of each sequence to each
fold or per-sequence predictions. Our full-record experiment samples its own
stratified 10-fold partition. It can compare descriptive means, but cannot
compute a paired uncertainty interval for their difference. The original
Veltri split, by contrast, is distributed verbatim and our test partition
matches ACEP's report; the published prediction vectors still are absent,
so even there a paired significance test cannot be reconstructed.

\subsection{Variance of the difference}
If $A$ and $B$ denote model scores across matched folds, the variance of
their mean difference would be
\begin{equation}
\mathrm{Var}(\overline{A}-\overline{B}) =
\mathrm{Var}(\overline{A}) + \mathrm{Var}(\overline{B})
- 2\mathrm{Cov}(\overline{A},\overline{B}).
\label{eq:pairedvariance}
\end{equation}
The covariance term depends on which cases each model classifies correctly
in the same folds. Published summary statistics do not supply that term.
Dropping it is not guaranteed conservative: positive covariance makes an
unpaired variance too large, negative covariance makes it too small.
Furthermore, ordinary fold scores share overlapping training records, so
simple independent-fold standard errors are at best approximate.

\subsection{Threshold and metric choices}
The reported accuracy and MCC use a fixed 0.5 threshold on an average of
uncalibrated sigmoid outputs. A monotone transformation of these outputs
preserves AUROC but can change the operating point. If we chose a threshold
after observing the test folds, accuracy and MCC would be optimistically
biased; therefore the full-record experiment never tunes a threshold on
the held-out fold. Comparisons to the published production model are still
not exact algorithmic replications: architecture, encoding, and preprocessing
are different by design. These differences are what is being tested, but
they also prevent attributing any score gap to a single component.

\subsection{Prevalence and precision}
The Feb2020 benchmark has equal positive and negative counts by
construction. The positive predictive value of a classifier with sensitivity
$Se$, specificity $Sp$, and target-population prevalence $\pi$ is
\begin{equation}
\mathrm{PPV}(\pi) = \frac{\pi Se}{\pi Se + (1-\pi)(1-Sp)}.
\label{eq:ppv}
\end{equation}
A high benchmark AUPRC at $\pi=1/2$ need not imply high precision when
screening an uncharacterized proteome where true AMP prevalence may be much
lower and is unknown. Our discovery scores therefore rank hypotheses; they
do not offer a validated per-candidate probability of activity.
